\section{Fundamentação Teórica}\label{sec:fundamentacao-teorica}

A fundamentação teórica pode ser subdividida em subseções de acordo com
o que for mais conveniente. A seguir são propostas formas de organização
sendo que a de correlatos é obrigatória. \textbf{Recomenda-se fortemente
que esta seção não ultrapasse quatro páginas}.

\subsection{Conceitos, técnicas e/ou ferramentas}\label{subsec:conceitos-tecnicas-ferramentas}

Na primeira parte deve-se abordar os \ul{conceitos, técnicas e/ou
ferramentas} mais relevantes envolvidos com o tema, devendo ser omitidas
metodologias de especificação e ferramentas de implementação que já são
conhecidas.

\subsection{Versão anterior do sistema/ferramenta/biblioteca/framework
(opcional)}\label{subsec:versao-anterior}

Quando o projeto propõe uma continuação ou extensão de um Trabalho de
Conclusão de Curso (TCC) anterior, deve-se descrevê-lo em uma seção
específica.

\subsection{Trabalhos correlatos}\label{subsec:trabalhos-correlatos}

Na segunda parte da fundamentação devem ser descritos os \ul{trabalhos
correlatos}. Devem ser incluídos preferencialmente trabalhos acadêmicos
com características e funcionalidades semelhantes ao que está sendo
produzido. A descrição deve ser feita em quadros conforme apresentado no
Quadro 1. \ul{Os itens apresentados no quadro são obrigatórios}. Quando
não for possível identificar, justificar no quadro.

\begin{center}\textbf{Quadro 1 -- Trabalho Correlato 1}\end{center}

\begin{longtable}[]{@{}
  >{\raggedright\arraybackslash}p{(\linewidth - 2\tabcolsep) * \real{0.1784}}
  >{\raggedright\arraybackslash}p{(\linewidth - 2\tabcolsep) * \real{0.8216}}@{}}
\toprule\noalign{}
\begin{minipage}[b]{\linewidth}\raggedright
Referência
\end{minipage} & \begin{minipage}[b]{\linewidth}\raggedright
\end{minipage} \\
\midrule\noalign{}
\endhead
\bottomrule\noalign{}
\endlastfoot
Objetivos & \\
Principais funcionalidades & \\
Ferramentas de desenvolvimento & \\
Resultados e conclusões & \\
\end{longtable}

Fonte: elaborado pelo autor.
